\documentclass[11pt]{article}
\usepackage[margin=1in]{geometry}
\usepackage{amsmath,amssymb,amsthm}
\usepackage[T1]{fontenc}
\usepackage{lmodern}
\usepackage{graphicx}
\usepackage[hidelinks]{hyperref}

\newtheorem{theorem}{Theorem}
\newtheorem{lemma}{Lemma}
\newtheorem{remark}{Remark}

\newcommand{\N}{\mathbb N}
\newcommand{\R}{\mathbb R}
\newcommand{\distmax}{\delta}

\title{The infinite-sum converse for statistically maximizing sequences fails}
\author{Codex agent lane 3 (GPT-5.6)}
\date{July 19, 2026}

\begin{document}
\maketitle

\section*{Status and Source Question}

\noindent\textbf{Status:} counterexample likely valid, pending human review.

Seal--Som--Basu--Dey study farthest-point maximizing sequences and
$M$-compactness in normed spaces \cite{SSBD}.  If $M$ is a nonempty bounded
subset of a normed space and $y$ is in the ambient space, write
\[
 \distmax(y,M)=\sup_{m\in M}\|y-m\|.
\]
A sequence $(x_n)$ in $M$ is \emph{statistically maximizing} when, for some
$y$, the scalar sequence $\|x_n-y\|$ converges statistically to
$\distmax(y,M)$.

Corollary 3.6 of the source says that statistical maximality in an infinite
$\ell_p$-sum implies statistical maximality in every coordinate.  Remark 3.7
asks whether the converse is true for $1\leq p<\infty$.

\begin{center}
\includegraphics[width=0.98\linewidth]{figures/open_problem_crop.png}
\end{center}

The answer is negative for every $1\leq p<\infty$.  The counterexample uses
only three-point subsets of two-dimensional Hilbert spaces.  In fact, every
coordinate sequence below is maximizing in the ordinary sense, while the
global sequence is not even statistically maximizing.

\section*{Definitions and Construction}

For a scalar sequence $(r_n)$, statistical convergence to $r$ means
\[
 \lim_{N\to\infty}\frac1N
 \bigl|\{n\leq N:|r_n-r|\geq\varepsilon\}\bigr|=0
 \qquad(\varepsilon>0).
\]
Fix $1\leq p<\infty$.  For $i\in\N$, set
\[
 \alpha_i=2^{-i},\qquad X_i=\R^2
\]
with the Euclidean norm, and define
\[
 A_i=\left(-\frac{\alpha_i}{2},0\right),\qquad
 B_i=\left( \frac{\alpha_i}{2},0\right),\qquad
 C_i=(\alpha_i,\alpha_i^2),
 \qquad M_i=\{A_i,B_i,C_i\}.
\]
Let
\[
 X=\left(\bigoplus_{i=1}^{\infty}X_i\right)_{\ell_p},
 \qquad
 M=\left(\bigoplus_{i=1}^{\infty}M_i\right)_{\ell_p}
   =\{(m_i)\in X:m_i\in M_i\text{ for every }i\}.
\]
This is a bounded nonempty subset of $X$: since $\alpha_i\leq1/2$,
\[
 \max_{m\in M_i}\|m\|_2
 =\|C_i\|_2
 =\alpha_i\sqrt{1+\alpha_i^2}
 \leq\frac{\sqrt5}{2}\alpha_i,
\]
and $(\alpha_i)\in\ell_p$.

Put $A=(A_i)_{i\geq1}$ and $B=(B_i)_{i\geq1}$, both elements of $M$, and
define
\[
 x_n=\begin{cases}
 A,&n\text{ even},\\
 B,&n\text{ odd}.
 \end{cases}
\]

\section*{Proof Intuition}

In coordinate $i$, the two alternating points $A_i,B_i$ are both farthest
from the harmless witness $(0,1)$.  Thus each coordinate sequence is as
strongly maximizing as possible: its distance to a witness is constantly
equal to the farthest distance.

The third point $C_i$ prevents these coordinate witnesses from being replaced
by a summable family.  If $A_i$ and $B_i$ are both farthest from a center
$(s_i,t_i)$, their equal distances force $s_i=0$.  Comparing either one with
$C_i$ then forces
\[
 t_i\geq\frac38+\frac{\alpha_i^2}{2}.
\]
Every simultaneous witness therefore has coordinate norm at least $3/8$.
Such a family cannot belong to an $\ell_p$-sum.

If the alternating global sequence were statistically maximizing, its even
and odd distances, each occurring on a set of density $1/2$, would both have
to equal the global farthest distance.  Coordinatewise additivity of the
$p$th power of the norm would then make every $A_i$ and $B_i$ simultaneously
farthest from the corresponding coordinate of one global witness.  The lower
bound above gives the contradiction.

\section*{Formal Counterexample}

\begin{lemma}[The coordinate obstruction]
\label{lem:coordinate}
For every $i$, the sequence alternating between $A_i$ and $B_i$ is an
ordinary maximizing sequence in $M_i$.  Moreover, if both $A_i$ and $B_i$
are farthest points of $M_i$ from $z_i=(s_i,t_i)$, then
\[
 s_i=0,
 \qquad
 t_i\geq\frac38+\frac{\alpha_i^2}{2},
 \qquad
 \|z_i\|_2\geq\frac38.
\]
\end{lemma}

\begin{proof}
Take $y_i=(0,1)$.  Direct calculation gives
\begin{align*}
 \|A_i-y_i\|_2^2=\|B_i-y_i\|_2^2
 &=1+\frac{\alpha_i^2}{4},\\
 \|C_i-y_i\|_2^2
 &=\alpha_i^2+(1-\alpha_i^2)^2
   =1-\alpha_i^2+\alpha_i^4.
\end{align*}
Their difference is
\[
 \alpha_i^2\left(\frac54-\alpha_i^2\right)>0
\]
because $0<\alpha_i\leq1/2$.  Hence $A_i$ and $B_i$ are both farthest from
$y_i$, and the alternating sequence has constant distance
$\distmax(y_i,M_i)$.

Now suppose that both $A_i$ and $B_i$ are farthest from
$z_i=(s_i,t_i)$.  They must be equidistant.  Since
\[
 \|z_i-A_i\|_2^2-\|z_i-B_i\|_2^2=2\alpha_i s_i,
\]
we obtain $s_i=0$.  Farthestness relative to $C_i$ gives
\begin{align*}
 t_i^2+\frac{\alpha_i^2}{4}
   =\|z_i-A_i\|_2^2
 &\geq \|z_i-C_i\|_2^2\\
 &=\alpha_i^2+(t_i-\alpha_i^2)^2.
\end{align*}
After cancellation and division by $2\alpha_i^2>0$, this is exactly
\[
 t_i\geq\frac38+\frac{\alpha_i^2}{2}.
\]
The norm bound follows.
\end{proof}

\begin{lemma}[Farthest distance in the bounded product]
\label{lem:product}
For every $z=(z_i)\in X$,
\[
 \distmax(z,M)^p
 =\sum_{i=1}^{\infty}\distmax(z_i,M_i)^p.
\]
If $u=(u_i)\in M$ is farthest from $z$, then each $u_i$ is farthest from
$z_i$ in $M_i$.
\end{lemma}

\begin{proof}
Every $M_i$ is finite, so choose $m_i\in M_i$ attaining
$\distmax(z_i,M_i)$.  The uniform coordinate bound used above implies
$(m_i)\in M$.  Therefore the coordinatewise upper bound for
$\|z-m\|_p^p$ is attained, proving the formula.

If $u$ is globally farthest, then
\[
 0=\distmax(z,M)^p-\|z-u\|_p^p
  =\sum_{i=1}^{\infty}
   \left(\distmax(z_i,M_i)^p-\|z_i-u_i\|_2^p\right).
\]
Every summand is nonnegative, so every summand is zero.
\end{proof}

\begin{theorem}[Negative answer to Remark 3.7]
\label{thm:counterexample}
For every $1\leq p<\infty$, there exist finite bounded sets
$M_i\subset\R^2$ and a sequence $(x_n)$ in their bounded $\ell_p$-sum such
that every coordinate sequence $(x_n^{(i)})$ is maximizing, hence
statistically maximizing, in $M_i$, but $(x_n)$ is not statistically
maximizing in the $\ell_p$-sum.
\end{theorem}

\begin{proof}
Use the construction above.  Lemma~\ref{lem:coordinate} proves that every
coordinate sequence is ordinary maximizing.

Suppose, toward a contradiction, that $(x_n)$ is statistically maximizing in
$M$.  There is then a $z\in X$ such that
\[
 \|x_n-z\|_p\longrightarrow_{\mathrm{stat}}\distmax(z,M).
\]
The scalar sequence on the left alternates between the two constants
$\|A-z\|_p$ and $\|B-z\|_p$.  The even and odd integers both have natural
density $1/2$.  Statistical convergence is therefore possible only if
\[
 \|A-z\|_p=\|B-z\|_p=\distmax(z,M).
\]
Thus both $A$ and $B$ are farthest points of $M$ from $z$.
Lemma~\ref{lem:product} implies that, for every $i$, both $A_i$ and $B_i$
are farthest from $z_i$ in $M_i$.  Lemma~\ref{lem:coordinate} now yields
\[
 \|z_i\|_2\geq\frac38\qquad(i\in\N).
\]
Consequently $\sum_i\|z_i\|_2^p=\infty$, contradicting $z\in X$.
\end{proof}

\section*{Verification Notes}

The proof is symbolic and uses no external theorem beyond the definitions in
the source.  The reusable script \texttt{code/check\_geometry.py} checks with
exact rational arithmetic, for $i=1,\dots,40$, both displayed coordinate
identities and the equality case at
$t_i=3/8+\alpha_i^2/2$.  It also prints finite-truncation lower bounds on a
putative witness norm for $p\in\{1,1.5,2,4,10\}$.  These calculations are a
sanity check, not part of the proof.

The most important review points are:
\begin{enumerate}
 \item the product formula in Lemma~\ref{lem:product}, including attainment
       by a coordinatewise selector in the bounded product $M$;
 \item the density-$1/2$ argument forcing both alternating global points to
       be farthest; and
 \item the sign and constant in the threshold
       $t_i\geq3/8+\alpha_i^2/2$.
\end{enumerate}

\section*{Novelty Audit and Limitations}

On July 19, 2026, the local run indexes were searched for arXiv:2008.08491,
$M$-compactness, statistical maximality, and infinite $\ell_p$-sums.  Web and
arXiv searches used the exact phrase from Remark 3.7, the exact source title
with citation/year variants, and close variants such as ``statistically
maximizing'' with ``direct sum'' or ``Banach.''  They returned the source
paper and source-adjacent earlier work on statistical convergence and
farthest points, but no later primary source answering this infinite-sum
converse.  This is a bounded novelty check, not an exhaustive citation audit.

The theorem answers exactly the statistical converse in Remark 3.7 and does
so simultaneously for every finite $p$.  It does not concern $p=\infty$, and
it makes no assertion about additional structural hypotheses on the sets
$M_i$ (for example, convexity) that might restore a converse.

\section*{Human-Review Recommendation}

Promote as a negative answer to Remark 3.7 if the three review points above
check out.  The construction is elementary, finite in every coordinate, and
the global obstruction is a direct nonsummability argument.

\begin{thebibliography}{9}
\bibitem{SSBD}
S. Seal, S. Som, S. Basu, and L. K. Dey,
\emph{Characterization of $M$-compact sets via statistically convergent
sequences}, arXiv:2008.08491v1, 2020.
\end{thebibliography}

\end{document}
